\documentclass[10pt,a4paper]{amsart}
\usepackage[margin=19mm]{geometry}
\usepackage{amsmath,amssymb,mathtools}
\usepackage[scr=rsfs,cal=euler]{mathalfa}
\usepackage{xfrac}
\usepackage[unicode,hidelinks,bookmarks=false]{hyperref}
\newcommand{\ie}{i.e.\ }
\newcommand{\cf}{cf.\ }
\newcommand{\resp}{respectively\ }
\newcommand{\Subsection}{\subsection}
\providecommand{\nobreakdash}{\nobreak\hbox{-}}
\setlength{\emergencystretch}{2em}
\multlinegap0pt
\usepackage{tikz}
\usetikzlibrary{cd}
\usepackage{amssymb}
\usepackage{mathtools}
\usepackage{tikz}
\usepackage{enumerate}
\usepackage[shortlabels]{enumitem}
\newtheorem{prop}{Proposition}
\newtheorem{lemma}{Lemma}
\newtheorem{thm}{Theorem}
\newcommand{\Bun}{{\mathrm{Bun}}}
\DeclareMathOperator{\Dmod}{D-mod}
\DeclareMathOperator{\IndCoh}{IndCoh}
\DeclareMathOperator{\LS}{LS}
\newcommand{\Nilp}{\mathrm{Nilp}}
\newcommand{\Poinc}{\operatorname{Poinc}}
\DeclareMathOperator{\QCoh}{QCoh}
\newcommand{\Whit}{\operatorname{Whit}}
\newcommand{\abc}{{\mathrm{abc}}}
\newcommand{\bL}{\mathbb{L}}
\newcommand{\spec}{\mathrm{spec}}
\title{The image of the constant sheaf under\\ the geometric Langlands equivalence}
\author{Kenta Suzuki}
\date{Source: arXiv:2608.25357v1 (CC BY 4.0)}
\begin{document}
\maketitle

\noindent\emph{Source and licence.} The image of the constant sheaf under\\ the geometric Langlands equivalence. Version arXiv:2608.25357v1 (CC BY 4.0), distributed by arXiv under the Creative Commons Attribution 4.0 International licence, http://creativecommons.org/licenses/by/4.0/, as recorded in the arXiv licence metadata for this exact version and shown on the abstract page https://arxiv.org/abs/2608.25357v1. Full licence notice: \texttt{licenses/arxiv-2608.25357-CC-BY-4.0.txt}.\par\smallskip

\noindent\textit{Editorial scope.} Counted unit: the article's thm conj:lafforgue-Z (source0.tex lines 745--754). The article's complete original proof of it is delivered with it (source0.tex lines 773--797, one \texttt{proof} environment of 25 lines, which the article places away from the statement). No mathematics is written, repaired or re-derived: the statement and the proof are the article's own lines, in the article's own notation, and no retained line is substituted or re-flowed.  Retained uncounted as the source-local context the proof needs: lines 296--306; lines 308--316; lines 559--561; lines 571--576; lines 761--770.  Nothing was omitted from any retained span, so the package carries no omission record.  No retained line contains a citation, so the wrapper carries no bibliography and no bibliography entry.  The retained lines contain the article's own diagram code, which the wrapper typesets with the standard tikz packages; no image file is delivered and the package declares no asset dependency.  Editorial formatting only, in addition: an AMS \texttt{amsart} preamble that restates the article's own declarations and macros used by the retained lines, namely the environments \texttt{prop}, \texttt{lemma}, \texttt{thm} on separate counters, together with the standard packages those lines need.  The retained environments print in this document as Proposition 1, Lemma 1, Theorem 1 in order of appearance, whereas the article prints the same items with the numbering of its own full document.  The complete native source of the article as distributed in the arXiv e-print for this exact version is bundled unchanged as \texttt{sources/208571/source0.tex}.
\begin{prop}
\label{lem:duality-intertwined}
    There is a commutative diagram
    \[
    \begin{tikzcd}
        \Dmod(\Bun_G)^\vee\arrow{r}{\mathbb D_{\Bun_G,!}}&\Dmod(\Bun_G)\arrow{dr}{\mathbb L_G}\\
        \IndCoh_\Nilp(\LS_{\check G})^\vee\arrow{r}{\mathbb D_{\mathrm{Serre}}}\arrow{u}{\mathbb L_G^\vee}&\IndCoh_\Nilp(\LS_{\check G})\arrow{r}{\tau[d]}&\IndCoh_\Nilp(\LS_{\check G}),
    \end{tikzcd}
    \]
    where $\tau$ denotes the Chevalley involution and $d=4\delta_G+2\delta_N$. 
\end{prop}

\begin{lemma}\label{lem:serre-duality}
    Let $\mathcal X$ and $\mathcal Y$ be QCA algebraic stacks. For any morphism $f\colon\mathcal X\to\mathcal Y$, there is a commutative diagram
    \[
    \begin{tikzcd}
        \IndCoh(\mathcal X)^\vee\arrow{r}{\mathbb D_{\mathrm{Serre}}}\arrow{d}{(f^!)^\vee}&\IndCoh(\mathcal X)\arrow{d}{f_*^{\IndCoh}}\\
        \IndCoh(\mathcal Y)^\vee\arrow{r}{\mathbb D_{\mathrm{Serre}}}&\IndCoh(\mathcal Y).
    \end{tikzcd}
    \]
\end{lemma}

\begin{lemma}\label{lem:dual-of-Z}
    Let $G$ be a reductive group, and let $\check Z$ be the center of the Langlands dual $\check G$. Then the dual of $\check Z$ is $G_\abc$.
\end{lemma}

\begin{prop}\label{prop:GLC-disconnected}
    Let $\check T$ be a possibly disconnected torus and let $T$ be the Langlands dual group. There is a canonical equivalence
    \begin{equation}\label{eq:disconnected-tori-equivalence}
    \mathbb L_T\colon \Dmod(\Bun_T)\simeq\QCoh(\LS_{\check T}).
    \end{equation}
\end{prop}

\begin{thm}\label{thm:laf1}
    There is a commutative diagram
    \[
    \begin{tikzcd}
        \Dmod(\Bun_G)\arrow{r}{\mathbb L_G}\arrow{d}{\pi_{\Bun_G!}}&\IndCoh_\Nilp(\LS_{\check G})\arrow{d}{\Whit^\spec[2(\delta_G+\delta_N-\delta_Z)]}\\
        \Dmod(\Bun_{G_{\mathrm{abc}}})\arrow{r}{\mathbb L_{G_{\mathrm{abc}}}}&\QCoh(\LS_{\check Z}),
    \end{tikzcd}
    \]
    where $G_{\mathrm{abc}}=G/[G,G]^{\mathrm{sc}}$ and $\bL_{G_\abc}$ is the equivalence of Proposition~\ref{prop:GLC-disconnected}.
\end{thm}

\begin{thm}\label{conj:lafforgue-Z}
    Let $\pi_{\Bun_G}\colon\Bun_G\to\Bun_{G_\abc}$ be the standard morphism. There is a commutative diagram
    \begin{equation}\label{eq:main-diagram}
    \begin{tikzcd}
    \Dmod(\Bun_{G_{\mathrm{abc}}})\arrow{r}{\mathbb L_{G_{\mathrm{abc}}}}\arrow{d}{\pi_{\Bun_G}^!}&\QCoh(\LS_{\check Z})\arrow{d}{\Poinc^\spec}\\
        \Dmod(\Bun_G)\arrow{r}{\mathbb L_G}&\IndCoh_\Nilp(\LS_{\check G})
    \end{tikzcd}
    \end{equation}
    where $G_{\mathrm{abc}}=G/[G,G]^{\mathrm{sc}}$ and $\bL_{G_\abc}$ is the equivalence of Proposition~\ref{prop:GLC-disconnected} (we use Lemma~\ref{lem:dual-of-Z} to identify the dual of $\check Z$ with $G_\abc$).
\end{thm}

\begin{proof}[Proof that Theorem~\ref{conj:lafforgue-Z} implies Theorem~\ref{thm:laf1}]
First of all, the miraculous duality functors for $G$ and $G_\abc$ sit in a commutative diagram
\begin{equation}\label{eq:miraculous-dual}
\begin{tikzcd}
\Dmod(\Bun_G)^\vee\arrow{r}{\mathbb D_!}\arrow{d}{(\pi^*_{\Bun_G})^\vee}&\Dmod(\Bun_G)\arrow{d}{\pi_{\Bun_G!}}\\
\Dmod(\Bun_{G_\abc})^\vee\arrow{r}{\mathbb D_!}&\Dmod(\Bun_{G_\abc})
\end{tikzcd}
\end{equation}
Indeed, by base change
\[
(\pi_{\Bun_G}\times1)^*\Delta_{\Bun_{G_\abc}!}e_{\Bun_{G_\abc}}=(1\times\pi_{\Bun_G})_!\Delta_{\Bun_G!}(e_{\Bun_G})\in\Dmod(\Bun_G\times\Bun_{G_\abc}),
\]
so for any $F\in\Dmod(\Bun_G)^\vee$,
\begin{equation}
(F\otimes\operatorname{Id})(\pi_{\Bun_G}\times1)^*\Delta_{\Bun_{G_\abc}!}(e_{\Bun_{G_\abc}})=(1\times\pi_{\Bun_G})_!(F\otimes\operatorname{Id})\Delta_{\Bun_G!}e_{\Bun_G}.
\end{equation}
Now there is a commutative diagram
\[
\begin{tikzcd}
    \Dmod(\Bun_G)\arrow[swap,bend left=10]{rrr}{\tau[2(\delta_G+\delta_N)]\circ\mathbb L_G}\arrow{d}{\pi_{\Bun_G!}}&\arrow{l}{\mathbb D_![2\delta_G]}\Dmod(\Bun_G)^\vee\arrow{d}{(\pi_{\Bun_G}^!)^\vee}&\arrow{l}{\mathbb L_G^\vee}\IndCoh_\Nilp(\LS_{\check G})^\vee\arrow{d}{(\Poinc^\spec)^\vee}\arrow[swap]{r}{\mathbb D_{\mathrm{Serre}}}&\IndCoh_\Nilp(\LS_{\check G})\arrow{d}{\Whit^\spec}\\
    \Dmod(\Bun_{G_{\mathrm{abc}}})\arrow[bend right=10]{rrr}{\tau[2\delta_Z]\circ\mathbb L_{G_\abc}}&\arrow[swap]{l}{\mathbb D_![\delta_Z]}\Dmod(\Bun_{G_{\mathrm{abc}}})^\vee\arrow{r}{\mathbb L_{G_{\mathrm{abc}}}^\vee}&\arrow[swap]{l}{\mathbb L_{G_\abc}^\vee}\QCoh(\LS_{\check Z})^\vee\arrow{r}{\mathbb D_{\mathrm{Serre}}}&\QCoh(\LS_{\check Z}).
    \end{tikzcd}
\]
Here, the left square is \eqref{eq:miraculous-dual}, the middle square is by assumption, the right square is by Lemma~\ref{lem:serre-duality}, and the top and bottom are by Proposition~\ref{lem:duality-intertwined}. The large commutative square is the desired statement.
\end{proof}

\end{document}
